% Szablon Inżynierskiego Artykułu Naukowego (PL)
% Autor: LaTeX Engineer
% ---------------------------------------------------------

\documentclass[11pt,a4paper,oneside]{article}

% --- Pakiety Podstawowe ---
\usepackage[T1]{fontenc}
\usepackage[utf8]{inputenc}
\usepackage[polish]{babel}
\usepackage{polski}
\usepackage{indentfirst}
\usepackage{geometry}
\geometry{
    a4paper,
    total={170mm,257mm},
    left=25mm,
    top=25mm,
}

% --- Typografia i Wygląd ---
\usepackage{microtype}
\usepackage{lmodern}
\usepackage{titlesec}
\usepackage{fancyhdr}
\usepackage{lastpage}

% --- Matematyka i Fizyka ---
\usepackage{amsmath, amsfonts, amssymb}
\usepackage{mathtools}
\usepackage{siunitx}
\sisetup{output-decimal-marker = {,}}

% --- Grafika i Tabele ---
\usepackage{graphicx}
\usepackage{booktabs}
\usepackage{caption}
\usepackage{subcaption}
\usepackage{float}

% --- Bibliografia i Linki ---
\usepackage[style=numeric, sorting=none, language=polish]{biblatex}
\addbibresource{references.bib}
\usepackage{hyperref}
\hypersetup{
    colorlinks=true,
    linkcolor=blue,
    filecolor=magenta,      
    urlcolor=cyan,
    citecolor=blue,
}

% --- Makra Inżynierskie ---
\newcommand{\unit}[1]{\,\mathrm{#1}}

% --- Nagłówek i Stopka ---
\pagestyle{fancy}
\fancyhf{}
\rhead{Inżynierski Artykuł Naukowy}
\lhead{\leftmark}
\rfoot{Strona \thepage \ z \pageref{LastPage}}

% --- Dane Dokumentu ---
\title{Profesjonalna Analiza Inżynierska Systemu X}
\author{Imię Nazwisko \\ \small Wydział Inżynierii, Politechnika ...}
\date{\today}

\begin{document}

\maketitle

\begin{abstract}
W niniejszym artykule przedstawiono kompleksową analizę inżynierską systemu X. Opisano metodykę pomiarową, wyniki eksperymentalne oraz wyciągnięto wnioski dotyczące optymalizacji procesu Y. Uzyskane wyniki wskazują na możliwość poprawy efektywności o \num{12.5}\%.
\end{abstract}

\section{Wstęp}
Współczesna inżynieria wymaga precyzyjnych narzędzi do dokumentowania wyników badań. LaTeX stanowi standard w~publikacjach technicznych, zapewniając najwyższą jakość typograficzną oraz spójność dokumentów.

W~sekcji \ref{sec:metodyka} opisano stanowisko badawcze, natomiast w~sekcji \ref{sec:wyniki} przedstawiono zebrane dane.

\section{Metodyka Badawcza}
\label{sec:metodyka}
Stanowisko badawcze składało się z~układu X oraz sensora Y. Pomiary przeprowadzono w~temperaturze pokojowej.

\begin{equation}
    P = V \cdot I \cdot \cos(\varphi)
\end{equation}

Gdzie:
\begin{itemize}
    \item $P$ -- moc czynna [W],
    \item $V$ -- napięcie [V],
    \item $I$ -- natężenie prądu [A],
    \item $\varphi$ -- kąt przesunięcia fazowego.
\end{itemize}

\section{Wyniki Eksperymentalne}
\label{sec:wyniki}
W tabeli \ref{tab:wyniki} zestawiono wyniki dla różnych wariantów obciążenia.

\begin{table}[H]
    \centering
    \caption{Porównanie parametrów pracy systemu.}
    \label{tab:wyniki}
    \begin{tabular}{lrrr}
        \toprule
        Wariant & Napięcie [V] & Prąd [A] & Sprawność (\%) \\
        \midrule
        Wariant A & \num{230.1} & \num{1.5} & \num{88.2} \\
        Wariant B & \num{231.5} & \num{2.1} & \num{91.5} \\
        Wariant C & \num{229.8} & \num{3.4} & \num{94.7} \\
        \bottomrule
    \end{tabular}
\end{table}

\section{Wnioski}
Analiza danych potwierdza hipotezę o~wzroście sprawności przy wzroście obciążenia w~zbadanym zakresie. Dalsze badania powinny skupić się na wpływie temperatury na stabilność pracy układu.

\printbibliography

\end{document}
